\section{Holonomía afín y simetrías de la memoria}
\label{nuclear:6}
\subsection{Composición efectiva de la memoria}

Para una emisión admisible \(\epsilon\), la construcción de \cref{iv:cont:concatenacion} define

\begin{equation}
\mathsf U_\epsilon(m)=C_\epsilon m+\kappa_\epsilon,
\qquad
\kappa_{\epsilon_2\epsilon_1}
=\kappa_{\epsilon_2}+C_{\epsilon_2}\kappa_{\epsilon_1}.
\label{nuc06:eq:1}
\end{equation}

El dominio es el módulo de memoria de la fuente; el codominio, el de la emisión siguiente. El producto conserva el orden. En una ruta \(\gamma\) de longitud \(r\) se obtiene

\begin{equation}
C_\gamma=C_r\cdots C_1,\qquad
\kappa_\gamma=\sum_{j=1}^r C_r\cdots C_{j+1}\kappa_j,
\qquad
m_\gamma=C_\gamma m_0+\kappa_\gamma.
\label{nuc06:eq:2}
\end{equation}

El producto vacío de la suma es la identidad. La demostración es la inducción de \eqref{nuc06:eq:1}: añadir la emisión siguiente multiplica el registro anterior por su transporte lineal y suma su incremento. La frontera posicional puede cancelarse mientras permanece la suma transportada de \eqref{nuc06:eq:2}.

\begin{proposition}[Clase de memoria de un retorno]

Sea \(M\) un módulo abeliano y sea una ruta cerrada en la base cuyo transporte de memoria sea \(H(m)=Cm+\kappa\), con \(C\) un automorfismo de \(M\). Su clase de retorno es

\begin{equation}
\mathfrak m(H)=[\kappa]\in M/(I-C)M.
\label{nuc06:eq:3}
\end{equation}

Un cambio de coordenadas de memoria \(m'=Rm+b\), con \(R\) invertible, transforma el retorno en

\begin{equation}
C'=RCR^{-1},\qquad
\kappa'=R\kappa+(I-C')b.
\label{nuc06:eq:4}
\end{equation}

La aplicación inducida por \(R\) lleva \([\kappa]\) a \([\kappa']\). El transporte \(H\) tiene un punto fijo exactamente cuando su clase \eqref{nuc06:eq:3} es cero.
\end{proposition}

\begin{proof}
Sustituir \(m=R^{-1}(m'-b)\) en \(H\) da \eqref{nuc06:eq:4}. Como \(R(I-C)=(I-C')R\), \(R\) induce un isomorfismo entre los cocientes. El término \((I-C')b\) representa cero en el cociente destino. Finalmente, la ecuación de punto fijo \(Cm+\kappa=m\) equivale a \((I-C)m=\kappa\). Esto prueba la última afirmación.
\end{proof}

La fórmula \eqref{nuc06:eq:3} permite decidir si un cambio de origen elimina el incremento o sólo lo redistribuye entre coordenadas. En particular, cuando \(C=I\), la clase es el propio vector \(\kappa\). Una traslación no nula permanece no nula bajo cualquier cambio afín invertible.

\begin{proposition}[Simetría del retorno frente a cambio de representación]

Una transformación afín \(A(m)=Rm+b\) conmuta con \(H\) exactamente cuando

\begin{equation}
RC=CR,\qquad (I-C)b=(I-R)\kappa.
\label{nuc06:eq:5}
\end{equation}
\end{proposition}

\begin{proof}
Las partes lineales de \(AH\) y \(HA\) son \(RC\) y \(CR\); sus partes constantes son \(R\kappa+b\) y \(Cb+\kappa\). Igualarlas da \eqref{nuc06:eq:5}. Si sólo se transporta \(H\) mediante \eqref{nuc06:eq:4}, no se requiere conmutación: se obtiene la misma ley en otra representación. Para una simetría del retorno fijado, en cambio, se requieren ambas igualdades de \eqref{nuc06:eq:5}.
\end{proof}

Así se distinguen con precisión la covarianza de la ley y el estabilizador de una memoria concreta. Esta diferencia es pertinente para estudiar variaciones de simetría del estado sin modificar las leyes de composición.

\subsection{Aplicación al registro efectivo de la vuelta observable}

La sección~\ref{vii:sec:retorno-memoria-gravedad} define dos contadores ordenados: \(g\) cuenta inversiones de orientación y \(s\) cambios efectivos de hoja. Su cómputo canónico por bloque de 27 pasos es

\[
c_{27}=(1,18).
\]

La fase observable completa retorna después de cuatro de esos bloques. En el grupoide de rutas, el módulo de registro es \(M=\mathbb Z^2\), y

\begin{equation}
c_{108}=(4,72)=4v,\qquad v=(1,18),\qquad
H(g,s)=(g+4,s+72).
\label{nuc06:eq:6}
\end{equation}

La identidad \eqref{nuc06:eq:6} es un resultado del corpus que se utiliza aquí; el comprobador verifica sus consecuencias, no vuelve a contar las emisiones originales. El símbolo \(c_{108}\) designa este incremento de dos contadores. Es distinto del sello dodecafásico \(K\) y del cociclo central de la representación de Heisenberg.

\begin{theorem}[Invariantes completos del registro bajo vueltas completas]

Las órbitas de \(H\) y su inversa sobre \(\mathbb Z^2\) quedan clasificadas exactamente por

\begin{equation}
\boxed{\quad \mathcal I(g,s)=s-18g,\qquad \nu(g,s)=[g]_4.\quad}
\label{nuc06:eq:7}
\end{equation}

En consecuencia,

\begin{equation}
\mathbb Z^2/\mathbb Z(4,72)\cong\mathbb Z\oplus\mathbb Z/4\mathbb Z.
\label{nuc06:eq:8}
\end{equation}
\end{theorem}

\begin{proof}
Añadir \((4,72)\) conserva ambas expresiones de \eqref{nuc06:eq:7}. Si dos registros tienen iguales invariantes, \(g'-g=4N\) para un entero \(N\) y \(s'-s=18(g'-g)=72N\); luego difieren exactamente en \(N(4,72)\). La aplicación de \eqref{nuc06:eq:7} es sobreyectiva: dados \(j\) entero y \(r\) módulo cuatro, elegir un representante \(g\) de \(r\) y poner \(s=j+18g\). Esto demuestra \eqref{nuc06:eq:8}.
\end{proof}

La descripción cubre todas las órbitas de este lector de memoria. Las órbitas del estado enriquecido completo conservan, además, sus restantes coordenadas. Para una prolongación dirigida hacia adelante, un registro posterior se obtiene cuando el entero N anterior es no negativo; la equivalencia bilateral de \eqref{nuc06:eq:8} pertenece al grupoide.

La magnitud \(\mathcal I\) es la única carga lineal entera primitiva, salvo signo, invariante bajo esas traslaciones. En efecto, para \(L(g,s)=ag+bs\), la invariancia exige \(4a+72b=0\), por lo que \((a,b)=b(-18,1)\). Una vuelta acumula memoria longitudinal y conserva una relación transversal determinada por el incremento.

El factor cuatro tiene aquí un origen preciso: es \(\gcd(4,72)\) y, en el cómputo anterior, el número de bloques de 27 que completa el retorno observable. La clase \(\nu\) conserva la información que se pierde al dividir el incremento por cuatro y guardar únicamente la dirección primitiva \(v\).

\begin{corollary}[Registro conjunto de bloque y fase]

Escríbase la fase de esos cuatro bloques como \(a\in\mathbb Z/4\mathbb Z\). La actualización

\[
(a,g,s)\longmapsto(a+1,g+1,s+18)
\]

conserva \((s-18g,[g-a]_4)\). Ambos datos clasifican sus órbitas bilaterales: la diferencia \(N=g'-g\) satisface a la vez \(a'-a=[N]_4\) y \(s'-s=18N\). Esto reúne el avance de memoria y el retorno de la fase sin reemplazar ninguno por el otro.
\end{corollary}

\subsection{Generador parabólico determinado por la memoria}

La pareja ordenada de contadores dota al módulo \(M\) de la forma alternante unimodular

\[
\Omega_M((g,s),(g',s'))=gs'-sg'.
\]

Su dominio es el módulo de registro. No se identifica esta forma con un área física dimensional, ni con la forma del toro APP, que tiene otro dominio. La elección de orden \((g,s)\) fija su orientación.

\begin{theorem}[Estabilizador unimodular del incremento]

Defínase, directamente desde el vector primitivo \(v\) de \eqref{nuc06:eq:6},

\begin{equation}
N_v m=v\,\Omega_M(v,m)=v\mathcal I(m),\qquad
N_v=\begin{pmatrix}-18&1\\-324&18\end{pmatrix}.
\label{nuc06:eq:9}
\end{equation}

Entonces \(N_v^2=0\), \(N_v\ne0\) y

\begin{equation}
\boxed{\quad
\operatorname{Stab}_{SL_2(\mathbb Z)}(c_{108})
=\{I+nN_v:n\in\mathbb Z\}.
\quad}
\label{nuc06:eq:10}
\end{equation}

Cada transformación de \eqref{nuc06:eq:10} conserva \(\Omega_M\), el incremento \(c_{108}\) y la carga \(\mathcal I\). Para \(n\) no nulo es una matriz unipotente no identidad, de traza dos, perteneciente al tipo parabólico.
\end{theorem}

\begin{proof}
\(\Omega_M(v,v)=0\) implica \(N_v^2m=v\Omega_M(v,v)\Omega_M(v,m)=0\). Su entrada superior derecha es uno, de modo que es no nulo. La matriz

\[
B=\begin{pmatrix}1&0\\18&1\end{pmatrix}
\]

tiene determinante uno, lleva el primer vector de base a \(v\) y satisface

\begin{equation}
N_v=BNB^{-1},\qquad N=\begin{pmatrix}0&1\\0&0\end{pmatrix}.
\label{nuc06:eq:11}
\end{equation}

Un automorfismo entero fija \(4v\) exactamente cuando fija \(v\), porque \(M\) carece de torsión. Una matriz de determinante uno que fija el primer vector de base tiene necesariamente la forma \(\begin{pmatrix}1&n\\0&1\end{pmatrix}\). Conjugarla por \(B\) prueba \eqref{nuc06:eq:10}. El resto se deduce de esa expresión o directamente de \eqref{nuc06:eq:9}.
\end{proof}

La fórmula \eqref{nuc06:eq:9} hace que la construcción sea independiente del segundo vector elegido para completar \(v\) a una base orientada unimodular: cualquier otro complemento difiere en un múltiplo entero de \(v\) y produce el mismo \(N_v\). Bajo una transformación orientada \(R\in SL_2(\mathbb Z)\), se tiene \(N_{Rv}=RN_vR^{-1}\).

La acción puede escribirse con especial claridad:

\begin{equation}
(I+nN_v)m=m+n\mathcal I(m)v.
\label{nuc06:eq:12}
\end{equation}

La carga transversal determina cuánto desplaza esa simetría a lo largo de la dirección de memoria. La suma de vueltas produce el desplazamiento \(4Nv\); las simetrías estabilizadoras preservan esa ley de avance.

Sobre el cociente \eqref{nuc06:eq:8}, la transformación \(I+nN_v\) actúa como \((\mathcal I,[g]_4)\mapsto(\mathcal I,[g+n\mathcal I]_4)\). Por tanto, conserva la colección de órbitas y puede permutar sus clases residuales; \(\nu\) es invariante del retorno \(H\), no de cada transformación estabilizadora. Esta diferencia se deduce directamente de \eqref{nuc06:eq:12}.

La conjugación \eqref{nuc06:eq:11} proporciona un mapa algebraico explícito desde el modelo parabólico del TRIT, definido en \cref{euler-trit-generadores}, al módulo de registro extendido a escalares reales. Se obtiene además \(\exp(tN_v)=I+tN_v\) y \(B\exp(tN)=\exp(tN_v)B\). Este entrelazamiento concierne a operadores: el movimiento particular seleccionado por el TPK utiliza también sus emisiones y calendario. El grupo aritmético completo de \eqref{nuc06:eq:10} fija este lector; el grupo de simetrías admisibles del estado total se obtiene intersectándolo con las acciones que preservan los demás datos y relaciones.

En el sistema trasladado \eqref{nuc06:eq:6}, los cambios afines \(m\mapsto Rm+b\) con \(R\) en \eqref{nuc06:eq:10} y \(b\) arbitrario conmutan con \(H\). Los que preservan además el registro inicial fijado deben satisfacer \(Rm_0+b=m_0\). Distinguir ambas condiciones separa simetría de la ley y simetría de una trayectoria señalada.

\begin{theorem}[Realización variacional de la carga y del estabilizador]
\label{nuc06:variacional}

La extensión escalar del registro, \(M_{\mathbb R}=M\otimes_{\mathbb Z}\mathbb R\), conserva las coordenadas generadas \((g,s)\) y la forma \(\Omega_M=dg\wedge ds\). Es una realización posterior del módulo entero. Fijando la convención \(\iota_{X_H}\Omega_M=dH\), se tiene

\begin{equation}
X_{\mathcal I}=v,\qquad
X_{\frac12\mathcal I^2}=\mathcal I v=N_vm.
\label{nuc06:eq:12a}
\end{equation}

Así, la traslación y la colección parabólica construidas admiten generadores hamiltonianos explícitos. El flujo de \(\mathcal I\) es \(m\mapsto m+uv\); en \(u\) igual a cuatro reproduce \eqref{nuc06:eq:6}. El flujo de \(\mathcal I^2/2\) es \(m\mapsto(I+uN_v)m\). El parámetro \(u\) corresponde a esta realización continua; el paso discreto de retorno y las emisiones intermedias mantienen su calendario TPK propio.
\end{theorem}

\begin{proof}
Para un vector \(X=(X_g,X_s)\), \(\iota_X(dg\wedge ds)=X_g\,ds-X_s\,dg\). Como \(d\mathcal I=ds-18\,dg\), el primer campo es \((1,18)=v\). La regla de la cadena da el segundo. \(\mathcal I\) permanece constante a lo largo de ambos campos, porque \(d\mathcal I(v)=0\). Por tanto el segundo flujo integra \(dm/du=\mathcal I(m_0)v\), equivalente a \(\exp(uN_v)=I+uN_v\). La no degeneración de \(\Omega_M\) implica que el Hamiltoniano que genera \(N_v\) es \(\mathcal I^2/2\) más una constante aditiva: su diferencial queda fijado por \(\iota_{N_vm}\Omega_M\).
\end{proof}

La acción matemática

\begin{equation}
\mathscr A[m]=\int_{u_0}^{u_1}
\left[s\frac{dg}{du}-\frac12(s-18g)^2\right]du
\label{nuc06:eq:12b}
\end{equation}

produce por variación las ecuaciones \(dg/du=\mathcal I\), \(ds/du=18\mathcal I\). La traslación \(g\mapsto g+\varepsilon\), \(s\mapsto s+18\varepsilon\) conserva el Hamiltoniano y cambia el lagrangiano en la derivada total \(d(18\varepsilon g)/du\), para \(\varepsilon\) constante. Su carga de Noether es

\begin{equation}
\underbrace{s\,\delta g}_{s}
-\underbrace{18g}_{\text{término de borde}}
=s-18g=\mathcal I.
\label{nuc06:eq:12c}
\end{equation}

La variación se interpreta por unidad de \(\varepsilon\). Las ecuaciones comprueban directamente \(d\mathcal I/du=18\mathcal I-18\mathcal I=0\). La propia traslación discreta \eqref{nuc06:eq:6} se realiza también mediante el Hamiltoniano \(H_{\rm ret}=4\mathcal I\), cuyo flujo de parámetro uno es \(H\); ambos Hamiltonianos conmutan en el sentido de Poisson. El flujo estabilizador \(I+uN_v\) y el retorno \(m+4v\) son acciones distintas que conmutan: en la recta \(\mathcal I=0\) el primero fija cada punto, mientras el segundo sigue acumulando memoria.

Más generalmente, todo Hamiltoniano \(C^1\) sobre este plano que sea invariante bajo las traslaciones \(m\mapsto m+uv\) tiene la forma \(F(\mathcal I)\). Para probarlo, las coordenadas globales \((g,\mathcal I)\) escriben esa simetría como traslación de \(g\); la invariancia hace independiente de \(g\) a \(H\). Su campo es \(F'(\mathcal I)v\). Esta clasificación explicita tanto la conservación común como la elección adicional de una dinámica concreta.

La acción \eqref{nuc06:eq:12b} es adimensional y pertenece al plano de registro realizado. Su demostración proporciona una relación variacional exacta entre memoria, carga y estabilizador. La conversión en acción física y en un tiempo dimensional se realiza mediante los lectores de acción y reloj correspondientes; no se asignan esas unidades al contador u.

\subsection{Curvatura de APP y conservación de la fase central}

La sección~\ref{iv:extension-curvatura-mixta} trabaja sobre \(X_9=(\mathbb Z/9\mathbb Z)^2\). Las hojas son \(S(x,y)=x+y\) y \(P(x,y)=xy\). Las conexiones mixtas utilizan diferencias de esos potenciales, no los potenciales como enlaces:

\begin{equation}
F(S,P)=\Delta_x\Delta_yP-\Delta_y\Delta_xS=1,
\qquad F(P,S)=-1\pmod9.
\label{nuc06:eq:13}
\end{equation}

La suma orientada de enlaces en el contorno de un rectángulo es \(ab\) o \(-ab\) según el orden. El área alternante es \(\sigma(u,v)=ad-bc\) para \(u=(a,b)\), \(v=(c,d)\). La sección~\ref{iv:weyl-area-extension} representa esos transportes mediante \(D(u)=X^aZ^b\), \(ZX=\omega XZ\), con \(\omega\) la raíz novena en la representación posterior. Entonces

\begin{equation}
D(u)D(v)=\omega^{bc}D(u+v),\qquad
D(u)D(v)D(u)^{-1}D(v)^{-1}=\omega^{-\sigma(u,v)}I.
\label{nuc06:eq:14}
\end{equation}

La posición resultante del conmutador retorna. La fase central retiene el área orientada con el signo fijado en \eqref{nuc06:eq:14}. Es una retención módulo nueve: un área múltiplo de nueve tiene fase central trivial, aunque un levantamiento con cociente y registro pueda distinguir sus recorridos. La fase no reemplaza ese levantamiento.

\begin{proposition}[Covarianza exacta del transporte central]

Como dos es invertible módulo nueve, defínase

\begin{equation}
W(a,b)=\omega^{ab/2}D(a,b),\qquad
W(u)W(v)=\omega^{-\sigma(u,v)/2}W(u+v).
\label{nuc06:eq:15}
\end{equation}

Para \(R\in SL_2(\mathbb Z/9\mathbb Z)\), la aplicación

\begin{equation}
\omega^zW(u)\longmapsto\omega^zW(Ru)
\label{nuc06:eq:16}
\end{equation}

es un automorfismo de la extensión central. Conserva sus productos y sus fases de conmutación.
\end{proposition}

\begin{proof}
La fase al multiplicar los dos \(W\) es
\(\frac{ab}{2}+\frac{cd}{2}+bc-\frac{(a+c)(b+d)}2=\frac{bc-ad}{2}\). Esto prueba \eqref{nuc06:eq:15}. La identidad \(\sigma(Ru,Rv)=\det(R)\sigma(u,v)=\sigma(u,v)\) muestra que \eqref{nuc06:eq:16} conserva la ley; \(R\) inversible proporciona el automorfismo inverso. Si \(\det R=-1\), el automorfismo correspondiente cambia también \(z\mapsto-z\); el signo de área y el centro se invierten conjuntamente.
\end{proof}

El registro íntegro y su fase central tienen tipos diferentes. La covarianza permite conservar la ley de área mientras cambia la presentación; la clase afín de \eqref{nuc06:eq:3} permite conservar el contenido de memoria mientras cambia su origen. Ambas propiedades actúan sobre lectores tipados de rutas, y sus mapas deben mantenerse separados al componerlos con la lectura excepcional del estado completo.

\begin{proposition}[Persistencia bajo marcos locales y subdivisiones]

Sea un lazo \(\gamma\) basado en \(x\) y sean sus enlaces invertibles \(U_\epsilon\). El cambio local \(U'_\epsilon=G_tU_\epsilon G_s^{-1}\) transforma su holonomía en \(H'_\gamma=G_xH_\gamma G_x^{-1}\): los marcos de los vértices interiores se cancelan sucesivamente. Si \(H_\gamma=\omega^kI\), la fase central permanece exactamente igual. Para enlaces refinados cuyo producto reproduce el enlace anterior, la subdivisión conserva asimismo \(H_\gamma\), por asociatividad. En las plaquetas APP la cancelación de aristas interiores materializa esta propiedad mediante Stokes finito.
\end{proposition}

Un rectángulo elemental tiene Wilson aditivo igual a uno y fase central no trivial. Por tanto, ese transporte mixto distingue una holonomía real de la identidad que producirían enlaces de la forma \(G_tG_s^{-1}\). La conservación bajo subdivisión concierne a lazos y enlaces compatibles: añadir otra área cambia el transporte.

El levantamiento entero de la holonomía distingue el área orientada de su residuo nonádico: para un rectángulo de lados cuatro y siete, \(28=1+9\cdot3\), mientras que \(-28=8+9(-4)\) corresponde a su orientación inversa. La fase residual coincide para áreas uno y veintiocho, mientras sus cocientes distinguen las acumulaciones. Este ejemplo exhibe la diferencia entre fase finita y registro íntegro; su cociente de área y el contador temporal \(q_{\rm mem}\) conservan dominios distintos.

\subsection{Composición entre los tres regímenes y continuidad del balance}

Las realizaciones de \cref{euler-trit-generadores,euler-trit-evaluaciones} construyen generadores \(J_k^2=-\tau_kI\) sobre sus fibras propias, de dimensión real dos. Se toman enlaces efectivamente declarados \(B_k:V_k\to V_{k+1}\) que preservan las formas alternantes, \(B_k^*\Omega_{k+1}B_k=\Omega_k\). El signo \(*\) en esta sección significa transposición/pullback bilineal. Los operadores \(E_k=\exp(t_kJ_k)\) preservan \(\Omega_k\): en dimensión dos, la traza nula del generador implica \(J_k^*\Omega_k+\Omega_kJ_k=0\). La identidad simpléctica es la hipótesis que se conservaría al extender el lema a dimensiones mayores.

Para la compresión nonádica y su complemento registrado,

\[
T_k=B_k(8I+E_k)/9,\qquad
D_k=B_k\sqrt8(E_k-I)/9,
\]

la expansión algebraica da

\begin{equation}
\Omega_k=T_k^*\Omega_{k+1}T_k+D_k^*\Omega_{k+1}D_k.
\label{nuc06:eq:17}
\end{equation}

En efecto, los términos cruzados se cancelan y el numerador es \(72\Omega_k+9E_k^*\Omega_kE_k=81\Omega_k\). Para \(P_0=I\), \(P_{k+1}=T_kP_k\), se obtiene, por sustitución sucesiva y en el orden de la ruta,

\begin{equation}
\boxed{\quad
\Omega_0=P_N^*\Omega_NP_N+
\sum_{k=0}^{N-1}(D_kP_k)^*\Omega_{k+1}(D_kP_k).
\quad}
\label{nuc06:eq:18}
\end{equation}

El terminal y cada complemento permanecen presentes. La identidad es bilineal orientada; los resultados de positividad se aplican en los perfiles positivos demostrados en las secciones anteriores. La forma conservada no obliga a que un enlace entre regímenes distintos entrelace sus generadores: si \(BJ_\tau=J_{\tau'}B\), elevar al cuadrado con \(B\) invertible impondría \(\tau=\tau'\). La conservación de una estructura común permite mantener diferenciados los tres regímenes.

El transporte íntegro \(U_N=(B_{N-1}E_{N-1})\cdots(B_0E_0)\) lleva \(V_0\) a \(V_N\). Una monodromía matricial requiere un cierre tipado \(K:V_N\to V_0\); sólo entonces \(KU_N\) es un endomorfismo cuya traza puede calcularse. Cambios de base compatibles conservan esa traza por conjugación. Retorno de fase, cierre de fibra y retorno del vector siguen siendo condiciones diferentes.

Si los enlaces tienen exclusivamente la forma \(B_{yx}=G_yG_x^{-1}\), su producto a lo largo de un camino es \(G_{\rm final}G_{\rm inicial}^{-1}\). Al cerrar con el mismo marco da identidad. La holonomía de \eqref{nuc06:eq:14} y el incremento de \eqref{nuc06:eq:6} muestran mecanismos concretos que esa mera sustitución de marcos no puede representar por sí sola.

\subsection{Compatibilidad con refinamiento y alcance arbitrario}

La acción de memoria es compatible con un homomorfismo de módulos de restricción \(p:M'\to M\) cuando

\begin{equation}
pC'=Cp,\qquad p\kappa'=\kappa.
\label{nuc06:eq:19}
\end{equation}

Entonces \(pH'=Hp\) y \(p\) induce un mapa de cocientes

\begin{equation}
M'/(I-C')M'\longrightarrow M/(I-C)M,
\qquad [\kappa']\longmapsto[\kappa].
\label{nuc06:eq:20}
\end{equation}

\begin{proof}
La primera igualdad de \eqref{nuc06:eq:19} lleva \((I-C')M'\) a \((I-C)M\); la segunda fija la imagen del incremento. La compatibilidad de \(H\) resulta por sustitución. Concatenaciones y refinamientos componen los mismos cuadrados gracias a la ley \eqref{nuc06:eq:1}.
\end{proof}

Una clase de retorno no nula que aparece en un nivel conservado permanece no nula en cualquier levantamiento compatible: una clase cero se proyectaría a cero. Esta afirmación se aplica al sistema de memorias y restricciones que satisface \eqref{nuc06:eq:19}; el constructor conjunto y los mapas de su sistema inverso se conservan en \cref{iv:continuo-conjunto}.

En el registro de \eqref{nuc06:eq:6}, para cada entero \(N\) se tiene \(H^N(m)=m+N(4,72)\). La prueba es algebraica para \(N\) arbitrario, no una inferencia desde una cantidad finita de iteraciones. En la prolongación hacia adelante, cada prefijo conserva su registro finito; la historia completa se describe mediante la colección compatible de esos registros, sin reemplazar sus contadores por una cantidad entera «infinita».

Si una dinámica completa tiene el lector \(q_{\rm mem}\) con \(q_{\rm mem}(\Gamma_9x)=q_{\rm mem}(x)+1\), tampoco puede poseer un estado periódico de período \(N\) positivo en ese dominio: aplicar el lector produciría \(q_{\rm mem}(x)=q_{\rm mem}(x)+N\). Este corolario utiliza el grado entero completo; sus cocientes residuales pueden tener períodos finitos, como corresponde al calendario observable del corpus.

\subsection{Consecuencia fundacional y alcance de la composición}

El enlace específico obtenido es

\begin{equation}
\text{incremento de retorno }4v
\ \longrightarrow\ \mathcal I(m)=\Omega_M(v,m)
\ \longrightarrow\ N_vm=v\mathcal I(m)
\ \longrightarrow\ I+nN_v.
\label{nuc06:eq:21}
\end{equation}

El registro acumulado determina una carga conservada y una colección parabólica que estabiliza su dirección. La evolución y la conservación pertenecen a la misma acción: cambian \(g\) y \(s\), permanece su relación \(s-18g\), y el incremento conserva un estabilizador explícito. Esto precisa una parte del vínculo entre memoria y simetría que motivó la investigación.

La igualdad \eqref{nuc06:eq:18} conecta esa investigación con los balances anteriores: una evolución puede redistribuir la forma entre lectura y complementos conservando el total. Las identidades de APP \eqref{nuc06:eq:13}–\eqref{nuc06:eq:16} muestran, en otro lector concreto del mismo origen, cómo el orden del transporte produce un dato central aun cuando retorna la posición.

La carga \eqref{nuc06:eq:7} es un invariante algebraico del transporte y, mediante la acción explícita \eqref{nuc06:eq:12b}, una carga de Noether de la realización del registro. El teorema~\ref{nuc06:variacional} y su desarrollo variacional construyen esa relación, su simetría y su término de borde. Los desarrollos anteriores de acción y corriente físicas conservan sus dominios y unidades; transportarlos a esta realización exige mantener sus mapas dimensionales.

La composición con los lectores dimensionales debe transportar el par \((v,\Omega_M)\) y la carga \(\mathcal I\), preservando sus unidades y la corriente variacional. Toda relación posterior con \(\pi,\phi,e,\alpha\) y \(K\) ha de proceder de esa composición; las fórmulas presentes fijan la información que debe conservarse.
